% 图契约：相邻端口对折后的外围网络；图中的端口、几何、箭头和方程对应本题。
% 版面：沿用前三模块的几何图+推导分区；目标印宽17cm，简图至中等密度。
% 布局：左侧真实结构，右侧关键关系；真实电路边不作为架构流程箭头处理。
\documentclass[tikz,border=7pt]{standalone}
\usepackage{fontspec,xeCJK,amsmath,amssymb}
\IfFontExistsTF{Times New Roman}{\setmainfont{Times New Roman}}{\setmainfont{TeX Gyre Termes}}
\IfFontExistsTF{Heiti SC}{\setCJKmainfont[AutoFakeBold=1.5]{Heiti SC}}{\setCJKmainfont[AutoFakeBold=1.5]{FandolHei-Regular}}
\usetikzlibrary{arrows.meta,bending,calc,patterns,positioning,decorations.pathmorphing}
\definecolor{ink}{HTML}{243745}
\definecolor{blue}{HTML}{275DAB}
\definecolor{green}{HTML}{177C63}
\definecolor{red}{HTML}{BD4935}
\definecolor{orange}{HTML}{AA6C16}
\definecolor{purple}{HTML}{715393}
\tikzset{every node/.style={font=\fontsize{9.3}{12.5}\selectfont,text=ink},
arrow/.style={-{Stealth[length=5pt 1.5,width'=0pt 3]},shorten >=2pt,shorten <=1pt,line width=1pt},
arrow short/.style={-{Stealth[length=3pt 1.0,width'=0pt 2.5]},shorten >=1pt,shorten <=0pt,line width=.8pt},
axis/.style={arrow,draw=ink!60,line width=.6pt},
vec/.style={arrow,draw=blue},
wire/.style={draw=ink!75,line width=.8pt},
guide/.style={dashed,draw=ink!45,line width=.5pt},
card/.style={rounded corners=2pt,draw=blue!65,fill=blue!4,align=center,text width=6.5cm,inner xsep=9pt,inner ysep=7pt},
title/.style={font=\fontsize{10.5}{14}\selectfont\bfseries,text width=15cm,align=center},
label/.style={font=\fontsize{8.4}{11}\selectfont,align=center},
note/.style={label,text width=6.1cm}}
\begin{document}
\begin{tikzpicture}[x=1cm,y=1cm]
\path[use as bounding box] (-8,-3.5) rectangle(8,3.65);
\node[title] at (0,3.25) {M4-053a\quad 相邻端口对折后的外围网络};
\begin{scope}[shift={(-4.05,0)}]
\draw[wire,draw=ink](-2.5,0)--(-1.975,0.0) (-1.4249999999999998,0.0)--(-0.9,0);\begin{scope}[shift={(-1.7,0.0)},rotate=0.0]\draw[draw=ink,fill=white,line width=.8pt](-0.275,-.085)rectangle(0.275,.085);\end{scope}
\node[label]at(-1.7,0.35999999999999993){$R/2$};
\draw[wire,draw=ink](-0.9,0)--(-0.275,0.0) (0.275,0.0)--(0.9,0);\begin{scope}[shift={(0.0,0.0)},rotate=0.0]\draw[draw=ink,fill=white,line width=.8pt](-0.275,-.085)rectangle(0.275,.085);\end{scope}
\node[label]at(0.0,0.36){$R/2$};
\draw[wire,draw=ink](0.9,0)--(1.4249999999999998,0.0) (1.975,0.0)--(2.5,0);\begin{scope}[shift={(1.7,0.0)},rotate=0.0]\draw[draw=ink,fill=white,line width=.8pt](-0.275,-.085)rectangle(0.275,.085);\end{scope}
\node[label]at(1.7000000000000002,0.35999999999999993){$R/2$};
\draw[wire,draw=ink](-0.9,0)--(-0.9,-0.475) (-0.9,-1.025)--(-0.9,-1.5);\begin{scope}[shift={(-0.9,-0.75)},rotate=-90.0]\draw[draw=ink,fill=white,line width=.8pt](-0.275,-.085)rectangle(0.275,.085);\end{scope}
\node[label]at(-0.54,-0.75){$R/2$};
\draw[wire,draw=ink](-0.9,-1.5)--(-0.275,-1.5) (0.275,-1.5)--(0.9,-1.5);\begin{scope}[shift={(0.0,-1.5)},rotate=0.0]\draw[draw=ink,fill=white,line width=.8pt](-0.275,-.085)rectangle(0.275,.085);\end{scope}
\node[label]at(0.0,-1.1400000000000001){$R$};
\draw[wire,draw=ink](0.9,-1.5)--(0.9,-1.025) (0.9,-0.475)--(0.9,0);\begin{scope}[shift={(0.9,-0.75)},rotate=90.0]\draw[draw=ink,fill=white,line width=.8pt](-0.275,-.085)rectangle(0.275,.085);\end{scope}
\node[label]at(0.54,-0.75){$R/2$};
\draw[wire](-2.5,0)--(-2.5,1.5)(2.5,0)--(2.5,1.5);\draw[wire,draw=ink](-2.5,1.5)--(-0.275,1.5) (0.275,1.5)--(2.5,1.5);\begin{scope}[shift={(0.0,1.5)},rotate=0.0]\draw[draw=ink,fill=white,line width=.8pt](-0.275,-.085)rectangle(0.275,.085);\end{scope}
\node[label]at(0.0,1.92){$R$};
\fill[ink] (-2.5,0)circle(2pt);\node[left=4pt] at(-2.5,0){$A$};
\fill[ink] (-0.9,0)circle(2pt);\node[above=4pt] at(-0.9,0){$X$};
\fill[ink] (0.9,0)circle(2pt);\node[above=4pt] at(0.9,0){$Y$};
\fill[ink] (2.5,0)circle(2pt);\node[right=4pt] at(2.5,0){$B$};
\fill[ink] (-0.9,-1.5)circle(2pt);\node[below=4pt] at(-0.9,-1.5){$H$};
\fill[ink] (0.9,-1.5)circle(2pt);\node[below=4pt] at(0.9,-1.5){$G$};
\node[note]at(0,-2.8){$X=\{D,E\}$，$Y=\{C,F\}$；上方为原直连棱 AB};\end{scope}
\node[card,draw=blue!65,fill=blue!4]at(3.9,1.8){$X=\{D,E\},\quad Y=\{C,F\}$};
\node[card,draw=green!65,fill=green!4]at(3.9,0){$R_{XY}=R/2\parallel2R$};
\node[card,draw=orange!65,fill=orange!4]at(3.9,-1.8){外围再与直连棱 $R$ 并联};
\end{tikzpicture}\end{document}